\subsection{Additional ablation of frequency jitter and nonzero-mode layers}
\label{sec:review-ablation}
We fixed two configurations,
$(m,K,q,\Delta t^{-1})=(37,72,66,128)$ and $(41,80,76,152)$,
before executing a two-by-two ablation: registered random paired-frequency
jitter $\xi_j\sim U[0.2,0.8]$ versus fixed midpoints $\xi_j=0.5$, and six
nonzero-mode analytic outer-layer features present versus absent.
The random cells use the five registered seeds (7331, 7687, 8011, 8353, 9041).
Each deterministic cell has one unique trajectory; unused seeds are not counted
as replicates. All cells retain the same 24 mean-mode functions, six mean-mode
layers, quadrature, and time step. Adapter checks confirm bitwise equality of
the mean-mode discretization arrays and implicit operators across the four cells.

We computed 16 new CPU trajectories to $T=72$: 14 ablation trajectories and
two enriched-random CPU controls. Ten frozen enriched-random GPU trajectories
were reused and independently rescored. Every trajectory was evaluated at the
same 19 physical checkpoints on three grids against both fine FD references,
using velocity and vorticity errors over the full domain and $r\leq2.5$.
All four errors must be at most 0.5\%.
Table~\ref{tab:review-ablation} gives maxima over seeds, times, grids, and references.
All 16 CPU runs completed without numerical failure; the 12 trajectories without
nonzero-mode layers failed the accuracy gate. The two CPU controls differ from
the archived GPU fields by less than $2.0\times10^{-10}$ in the worst field metric.

\begin{table}[htbp]
\centering\small
\caption{Additional CPU accuracy ablation. Errors are percentages. Pass counts
exclude the two separate CPU controls. Random rows report the worst of five
seeds; midpoint rows contain one deterministic configuration.}
\label{tab:review-ablation}
\begin{tabular}{lrrrrrr}
\hline
Configuration & Rank & $e_u$ & $e_\omega$ & $e_{u,\mathrm{near}}$ & $e_{\omega,\mathrm{near}}$ & Pass \\
\hline
Mid, Random+layer & 39 & 0.0294 & 0.4096 & 0.0251 & 0.2100 & 5/5 \\
Mid, Midpoint+layer & 39 & 0.0293 & 0.4092 & 0.0249 & 0.2097 & 1/1 \\
Mid, Random/no layer & 34 & 19.4054 & 73.8980 & 26.4854 & 48.6306 & 0/5 \\
Mid, Midpoint/no layer & 34 & 19.3644 & 73.7836 & 26.4397 & 48.5825 & 0/1 \\
High, Random+layer & 40 & 0.0182 & 0.3319 & 0.0148 & 0.1799 & 5/5 \\
High, Midpoint+layer & 40 & 0.0181 & 0.3312 & 0.0147 & 0.1794 & 1/1 \\
High, Random/no layer & 35 & 18.7540 & 72.0115 & 25.7262 & 47.7640 & 0/5 \\
High, Midpoint/no layer & 35 & 18.7060 & 71.8705 & 25.6678 & 47.6945 & 0/1 \\
\hline
\end{tabular}
\end{table}

Both midpoint-plus-layer configurations pass, with errors comparable to the
random enriched cells. These data therefore do not establish randomness as
necessary for the observed accuracy. With raw $m$ and the SVD cutoff fixed,
removing the nonzero-mode layers reduces retained rank from 39/40 to 34/35 and
produces large velocity and vorticity errors; the mean rank remains 30.
The compact configurations depend strongly on analytic outer-layer enrichment,
but this design does not separate enrichment geometry from the increased
effective dimension. It is not an equal-rank enrichment ablation, and only one
deterministic frequency design was tested. All four error maxima for both
deterministic no-layer configurations occur at $T=72$. Their maximum vorticity
errors for $t\leq16$ are 2.820\% and 2.592\%, growing to 73.784\% and 71.871\%
at later times. The large errors therefore are not solely a consequence of a
near-zero startup reference field; the mechanism of error amplification is
not established by these data. A subsequent construction-only check increased raw $m$ to
41, 45, 49, 53, 57, and 61 while retaining bandwidth 62, quadrature, and the
SVD cutoff. Both resolutions produced ranks 35, 35, 36, 36, 36, and 36,
respectively, never recovering rank 39 or 40. We therefore did not add a
purported equal-rank $T=72$ control or tune the bandwidth/cutoff after failure.
Concurrent CPU elapsed times are
diagnostic only: they provide neither isolated timing nor new GPU speedup
evidence, and do not extend the original 78 GPU qualification trajectories.
